다음 그림과 같이 점 $\mathrm{O}$에서 $\overline{\mathrm{AC}}$에 내린 수선의 발을 $\mathrm{F}$라고 하면 $\triangle\mathrm{OAD}$와 $\triangle\mathrm{OAF}$에서

$\angle\mathrm{ODA} = \angle\mathrm{OFA} = 90^\circ$,

$\overline{\mathrm{OA}}$는 공통, $\angle\mathrm{OAD} = \angle\mathrm{OAF}$

이므로 $\triangle\mathrm{OAD} \equiv \triangle\mathrm{OAF}$ (RHA 합동)

$\therefore \overline{\mathrm{AD}} = \overline{\mathrm{AF}}$, $\overline{\mathrm{OD}} = \overline{\mathrm{OF}}$ $\cdots\cdots\mbox{㉠}$

또, $\triangle\mathrm{OCF}$와 $\triangle\mathrm{OCE}$에서

$\angle\mathrm{OFC} = \angle\mathrm{OEC} = 90^\circ$,

$\overline{\mathrm{OC}}$는 공통, $\angle\mathrm{OCF} = \angle\mathrm{OCE}$

이므로 $\triangle\mathrm{OCF} \equiv \triangle\mathrm{OCE}$ (RHA 합동)

$\therefore \overline{\mathrm{CF}} = \overline{\mathrm{CE}}$, $\overline{\mathrm{OF}} = \overline{\mathrm{OE}}$ $\cdots\cdots\mbox{㉡}$

㉠, ㉡에 의하여 $\overline{\mathrm{OD}} = \overline{\mathrm{OE}}$

다음 그림과 같이 $\overline{\mathrm{OB}}$를 그으면 $\triangle\mathrm{ODB}$와 $\triangle\mathrm{OEB}$에서

$\angle\mathrm{ODB} = \angle\mathrm{OEB} = 90^\circ$, $\overline{\mathrm{OB}}$는 공통, $\overline{\mathrm{OD}} = \overline{\mathrm{OE}}$

이므로 $\triangle\mathrm{OBD} \equiv \triangle\mathrm{OBE}$ (RHS 합동)

$\therefore \overline{\mathrm{BD}} = \overline{\mathrm{BE}}$ $\cdots\cdots\mbox{㉢}$

$\overline{\mathrm{CF}} = \overline{\mathrm{CE}} = x\,\mathrm{cm}$라고 하면 $\overline{\mathrm{AD}} = \overline{\mathrm{AF}} = (5 - x)\,\mathrm{cm}$,

$\overline{\mathrm{BD}} = 7 + (5 - x) = 12 - x\,(\mathrm{cm})$, $\overline{\mathrm{BE}} = (10 + x)\,\mathrm{cm}$

이므로 ㉢에 의하여

$12 - x = 10 + x$, $2x = 2$ $\therefore x = 1$

따라서 $\overline{\mathrm{CE}}$의 길이는 $1\,\mathrm{cm}$이다.
